% Bogurodzica: working edition of stanzas 1–2 with melody
% Build:  gregorio bogurodzica.gabc   (writes bogurodzica.gtex)
%         lualatex bogurodzica-edition.tex   (run twice; needs GregorioTeX 6 and the Junicode font)
\documentclass[11pt,a4paper]{article}
\usepackage[margin=2.2cm,top=2cm,bottom=2.4cm]{geometry}
\usepackage{fontspec}
\setmainfont{Junicode}[Path=/usr/share/fonts/opentype/junicode/,Extension=.otf,UprightFont=*-Regular,ItalicFont=*-Italic,BoldFont=*-SmBold,BoldItalicFont=*-SmBoldItalic,Numbers=OldStyle]
\usepackage{polyglossia}
\setdefaultlanguage[variant=british]{english}
\setotherlanguage{polish}
\usepackage{microtype}
\usepackage{gregoriotex}
\usepackage{booktabs,tabularx,longtable,array}
\usepackage{enumitem}
\usepackage{multicol}
\usepackage[hidelinks]{hyperref}
\usepackage{xcolor}
\definecolor{rubric}{RGB}{160,30,30}

\setlength{\parindent}{0pt}
\setlength{\parskip}{0.55em}
\linespread{1.08}

% ---- Gregorio settings -----------------------------------------------------
\gresetinitiallines{1}
\grechangestyle{initial}{\fontsize{40}{40}\selectfont}
\grechangestaffsize{19}
\grechangedim{baselineskip}{60pt plus 5pt minus 5pt}{scalable}
\grechangestyle{annotation}{\small\color{rubric}}

\newcommand{\sig}[1]{\textup{\textsc{\MakeLowercase{#1}}}}          % siglum
\newcommand{\pl}[1]{\foreignlanguage{polish}{\textit{#1}}}
\newcommand{\dipl}[1]{\foreignlanguage{polish}{\textup{#1}}}
\newcommand{\lem}[1]{\foreignlanguage{polish}{#1}\,\textbf{]}}
\newcommand{\pn}[1]{\textsc{#1}}            % pitch string
\newcommand{\sect}[1]{\section*{#1}\addcontentsline{toc}{section}{#1}}

\title{\textsc{Bogurodzica}\\[0.3em]
  \large The two oldest stanzas and their melody\\
  \normalsize A critical edition in square notation (GregorioTeX), with apparatus}
\author{Edited from Kraków, Biblioteka Jagiellońska, MS 1619}
\date{Working edition, October 2026}

\begin{document}
\maketitle
\thispagestyle{empty}

\begin{center}\small
\begin{minipage}{0.86\textwidth}
\textit{Status.} Working critical edition. The melody follows the oldest notated
witness, Kraków, BJ 1619. I collated it neume by neume against the
high-resolution Polona facsimile and against Chybiński's reading (1907); they
agree in every pitch. Variants from three later musical witnesses come at
second hand, from Chybiński's comparative examples. Readings of the standard
edition (Woronczak, Ostrowska \& Feicht 1962) are cited from Marek Adamiec's
digital transcriptions based on it and from two 1962 reviews. Points that only
the printed 1962 volume can settle are flagged \textbf{[W1962]}.
\end{minipage}
\end{center}

% ============================================================================
\sect{Edition}

\input{bogurodzica.gtex}

\vspace{0.4em}
{\footnotesize Editorial bar lines (the source has no divisiones). Final
\textit{mora vocis} only at the end of each stanza. Neume grouping follows
\sig{K}: climacus on \pl{-dzi-} of \pl{Bogurodzica}, porrectus on \pl{-wi-},
three separate puncta on \pl{-dzi-na} of \pl{Gospodzina}.}

% ----------------------------------------------------------------------------
\newpage
\subsection*{Text}

\begin{multicols}{2}\small
\textbf{Diplomatic} (BJ 1619, after Łoś 1922; \dipl{ǫ} = nasal vowel)\\[0.3em]
\dipl{Bogv rodzicza dzewicza bogem slawena maria\\
Vtwego syna gospodzina matko swolena maria\\
Sziszczi nam spwczi nam kyrieleyson}\\[0.4em]
\dipl{Twego dzela krzcziczela boszicze\\
Vslisz glosi naplen misli czlowecze\\
Slisz modlitwǫ yǫsz nosimi\\
A dacz raczi gegosz prosimi\\
a naswecze zbozni pobith\\
posziwocze raski pzbith ky'eleyson}

\columnbreak
\textbf{Critical (normalised, after Woronczak)}\\[0.3em]
\pl{Bogurodzica dziewica, Bogiem sławiena Maryja,\\
u twego syna Gospodzina matko zwolena, Maryja!\\
Zyszczy nam, spuści nam. Kirielejson.}\\[0.4em]
\pl{Twego dziela Krzciciela, Bożyce,\\
usłysz głosy, napełń myśli człowiecze.\\
Słysz modlitwę, jąż nosimy,\\
a dać raczy, jegoż prosimy:\\
a na świecie zbożny pobyt,\\
po żywocie rajski przebyt. Kirielejson.}
\end{multicols}

{\footnotesize Pilat (1879) reads \sig{K} differently in four places:
\dipl{slavena}, \dipl{O dacz}, \dipl{prosimy}, \dipl{przebith}. The facsimile
supports Łoś at \dipl{prosimi} and \dipl{pzbith}. For \dipl{A/O dacz} see the
palaeographical notes. The score keeps the modern seven-syllable underlay
\pl{Ky-ri-e e-le-i-son}.\par}

\textbf{Translation.} \textit{Mother of God, Virgin glorified by God, Mary!
Chosen Mother, Mary, win for us from thy Son the Lord, send down to us. Kyrie
eleison.} \textit{For thy Baptist's sake, Son of God, hear our voices, fulfil
the thoughts of men. Hear the prayer we bring, and deign to grant what we ask:
on earth a godly life, after life a dwelling in paradise. Kyrie eleison.}

% ============================================================================
\sect{Sources and sigla}

\textit{Warning on sigla.} Older literature numbers the Kraków copies the
other way round. Pilat (1879) and Kalina (1880) call BJ 408 ``Kraków~I'' (it
was then dated 1408) and the notated BJ 1619 ``Kraków~II''. Łoś (1922) and
later scholars reverse this. To avoid the clash this edition uses
shelfmark-based sigla.

\subsection*{Musical witnesses}
\begin{description}[leftmargin=2.6em,style=nextline,itemsep=0.35em,font=\normalfont]
\item[\sig{K}] \emph{Kraków, Biblioteka Jagiellońska, MS 1619} (\emph{base
  text}). Added on the back flyleaf of a copy of Latin sermons made by Maciej of
  Grochów, vicar of Kcynia (hence the ``Kcynia copy''). Woronczak dates it
  before 1408; it is usually given as c.\,1407. Stanzas 1--2 only, with music
  throughout, in seven systems. Each musical section opens with a decorated
  initial: \pl{B(ogu)}, \pl{U (twego)}, \pl{Z(yszczy)}; \pl{T(wego)},
  \pl{U(słysz)}, \pl{S(łysz)}, \pl{A (dać)} (Kalina). Gothic chant notation of
  the German--Messine (``Hufnagel'') type, with C and F clefs and no custodes.
  For staff lines, clefs and doubtful notes see the palaeographical notes.
  Interlinear Latin in red in system~5 (\textit{virtualis ingressus \dots}) is
  a later addition.
\item[\sig{K70}] \emph{Kraków, Biblioteka Jagiellońska, manuscript dated
  1570} (shelfmark not given by Chybiński \textbf{[W1962]}). Late
  ``popularised'' form that tends towards B\,flat before cadences on~a;
  Chybiński counts seven cases.
\item[\sig{Gn}] \emph{Gniezno cathedral tradition}. Chybiński cites a Gniezno
  manuscript that agrees with \sig{K70} in the opening phrase. Kalina describes
  the ``Gniezno authentic'' melody, printed in \emph{Przyjaciel Ludu} (Leszno)
  4, no.~46 (1838), p.~364. It covers the whole text, uses smaller note values
  and sharps, and was ``still sung in Gniezno'' in 1880.
\item[\sig{W}] \emph{Warsaw, library of the Literati Chapel at St John's
  church, MS 998 D}. In mensural notation, so it is the only witness with
  explicit rhythm. Chybiński treats it as the best of the later sources. Its
  present whereabouts need checking \textbf{[W1962]}.
\item[\sig{Cz}] \emph{Częstochowa (Jasna Góra), manuscript of the late 15th
  century}. The manuscript is lost and survives in a drawing reproduced by
  Woronczak (online in Adamiec's edition). Five-line staves carry notation for
  the first two lines of stanza~1 only, with a short neumed passage further
  down. Kalina reports variants in its first five notes, at the end of the first
  half-period, and in some middle notes. It has no music for \pl{zyszczy nam,
  spuści nam}. Chybiński dismissed it.
\item[\sig{Z}] \emph{Warsaw, Zamoyski Library, MS 1152} (1672). Mensural.
  Chybiński dismissed it as badly written and noted that it transmits the
  same version as \sig{K70}, which writes it more correctly.
\item[\sig{L}] \emph{Latin translation \emph{Dei genitrix illibata}} (early 16th
  century per Woronczak; text in Adamiec). Its melody is edited by Feicht 1975,
  pp.~187, 189, and is cited here from Flotzinger.
\end{description}

\subsection*{Text witnesses collated for stanzas 1--2}
\begin{description}[leftmargin=2.6em,style=nextline,itemsep=0.25em,font=\normalfont]
\item[\sig{K408}] Kraków, BJ MS 408, f.~87\textsuperscript{v}, added in the
  early 15th century to a codex of 1408 (\emph{Decisiones Rotae}). Stanzas
  1--14, no music.
\item[\sig{Wa}] Warsaw, Biblioteka Narodowa, from Holy Cross (Łysiec), second
  half of the 15th century (c.\,1456). 19 stanzas.
\item[\sig{Cz}] Częstochowa, as above.
\item[\sig{K3}] Kraków, manuscript of the first half of the 16th century
  (Pilat's ``Kraków III''), with 23 stanzas.
\item[\sig{Ł}] Jan Łaski, \emph{Commune incliti Poloniae Regni privilegium}
  (Kraków: J.\,Haller, 1506), printed in the \emph{Registrum} as \emph{prima
  omnium devotissima \dots\ cantio}, attributed to St Adalbert. \emph{Text
  only}: the facsimile shows no notation.
\item[\sig{M}] Mateusz z Kościana, 1543 (Gniezno group).
\item[\sig{H}] Herburt, 1570. \quad \sig{S}~Skarga, \emph{Żywoty świętych}, 1579.
\end{description}
Readings of \sig{K408}, \sig{Wa}, \sig{Cz}, \sig{K3}, \sig{M}, \sig{H} and
\sig{S} are taken from Pilat's synoptic table (1879, table~7), checked against
Woronczak's normalised transcriptions where these exist.

% ============================================================================
\sect{Editorial principles}
\begin{enumerate}[itemsep=0.2em]
\item \textbf{Base text.} \sig{K} throughout. Three readings of the melody were
  compared: the facsimile, Chybiński's transcription, and the standard modern
  transcription (the version circulated with the tempo marking
  \textit{Uroczyście}). All three agree in every pitch.
\item \textbf{Pitch.} Transcribed at the pitch fixed by \sig{K}'s clefs: final
  D, ambitus C--d$'$, mode~1 (Chybiński). With a C clef on the fourth line, the
  gabc letters \texttt{c d e f g h i j k} equal C D E F G a b c$'$ d$'$. No
  flats are supplied. The only b (\pl{spuś-ci}) is natural, and the B-flats of
  \sig{K70} belong to a later layer.
\item \textbf{Rhythm.} \sig{K} has no rhythmic notation, so the score is
  unmeasured square notation. The mensural reading of \sig{W} and Chybiński's
  duple-metre reconstruction are reported but not adopted.
\item \textbf{Text.} A normalised form of \sig{K} following Woronczak's
  transcription. This edition keeps \pl{u} at the start of 1.2, as \sig{K}
  does and as its notes require. Woronczak's reconstructed text of the
  original song omits it. Note also \pl{dziela} (`for the sake of', not
  \pl{dzieła}), \pl{zyszczy} (not \pl{ziści}) and \pl{Bożyce}.
\item \textbf{Bar lines} are editorial and mark the verse structure.
\end{enumerate}

% ============================================================================
\sect{Musical apparatus}

{\small
Pitch names: C D E F G a b c$'$ d$'$; hyphens join notes sung on one syllable.
Readings of \sig{K70} and \sig{W} are taken from Chybiński's modern-notation
comparative examples (1907, pp.~30--41). Rhythm is left out and only pitch
content is given. ``(?!)'' marks places Chybiński himself flagged as
corrupt.\par}

\renewcommand{\arraystretch}{1.25}
\begin{longtable}{@{}p{0.9cm}>{\raggedright\arraybackslash}p{3.2cm}>{\raggedright\arraybackslash}p{11.6cm}@{}}
\toprule
\textbf{Loc.} & \textbf{Lemma} & \textbf{Readings}\\
\midrule\endhead
\multicolumn{3}{@{}l}{\textit{Stanza 1}}\\
1.1 & \pl{Bogurodzica} & \sig{K} D C F E-D-C D \quad \sig{K70}, \sig{Gn} D D-C F E-C-E D \quad \sig{W} D C F E D \quad \sig{L} has D-C where \sig{K} has E-D-C on \pl{-dzi-} (Flotzinger)\\
1.2 & \pl{dziewica} & \sig{K} F G-F-G a \quad \sig{K70} F G b\,flat a \quad \sig{W} F G a\\
1.3 & \pl{Bogiem sławiena} & \sig{K} c$'$ d$'$ c$'$ a-G-F F \quad \sig{K70} \textbf{b} d$'$ c$'$ a-G F (?!) \quad \sig{W} c$'$ d$'$ c$'$ a F\\
1.4 & \pl{Maryja} & \sig{K} F-G-a-G E-C D \quad \sig{K70} F-a-G-E C E D (?!) \quad \sig{W} F a G E D D (?!)\\
1.5 & \pl{U twego syna} & \sig{K} d$'$ c$'$ a-G-F G-F-G a \quad \sig{K70} \textit{om.} d$'$ c$'$; a G F G b\,flat a a (alternatively a G F F G b\,flat a) \quad \sig{W} \textit{om.} d$'$ c$'$; simplified four-note form beginning on a \quad \sig{L} a-a on \textit{Apud} (Flotzinger, who suggests the same repeated a for \pl{U twe-})\\
1.6 & \pl{Gospodzina} & \sig{K} c$'$ d$'$ c$'$\,a\,G F \quad \sig{K70} c$'$ d$'$ c$'$ a-G F \quad \sig{W} c$'$ d$'$ c$'$ a F\\
1.7 & \pl{Matko} & \sig{K} c$'$ a-G-F \quad \sig{K70} c$'$ a-G (?!) \quad \sig{W} c$'$ a\\
1.8 & \pl{zwolena, Maryja} & as 1.4 in all witnesses (Chybiński)\\
1.9 & \pl{Zyszczy nam} & \sig{K} F G-F-G a \quad \sig{K70} F G b\,flat a \quad \sig{W} F G a\\
1.10 & \pl{spuści nam} & \sig{K} c$'$ b\,natural d$'$ \quad \sig{K70} c$'$ b\,natural d$'$ d$'$ (?!) \quad \sig{W} = \sig{K}\\
1.11 & \pl{Kyrie} & \sig{K} c$'$ a G-F \quad \sig{K70}, \sig{W} c$'$ a F\\
\midrule
\multicolumn{3}{@{}l}{\textit{Stanza 2}}\\
2.1 & \pl{Twego dziela Krzciciela} & \sig{K} D-E D-C F G E C D \quad \sig{K70} D D-C F F G G E E C D D (?!; ``agrees with \sig{K} but badly written'') \quad \sig{W} D D-C F G E C-E D\\
2.2 & \pl{Bożyce} & as 1.4 in all witnesses (``repeated for the fifth time'')\\
2.3 & \pl{usłysz głosy} & \sig{K} a G-F G-F-G a \quad \sig{K70} a G-F G b\,flat a \quad \sig{W} a G-F G a a\\
2.4 & \pl{napełń myśli człowiecze} & \sig{K} D C F G E C D \quad \sig{K70}, \sig{W} read \pl{napełnij}, with a consequent rhythmic change; Chybiński gives no pitches\\
2.5 & \pl{Słysz modlitwę, jąż nosimy} & \sig{K} C F G E-D-E C D E D \quad \sig{K70} \emph{different melody}, split into two halves: C G a b\,flat a a \,\textbar\, D D C C F F E D \quad \sig{W} \emph{different}: C F G a (fermata) \,\textbar\, D D C F G a a, ending on the confinalis a (rejected by Chybiński)\\
2.6 & \pl{a dać raczy} & \sig{K} a G-F G-F-G a \quad \sig{K70} a G F G b\,flat a a (?!)\\
2.7 & \pl{jegoż prosimy} & \sig{K} D-E D-C F G a \quad \sig{K70} D D-C F G b\,flat a\\
2.8 & \pl{a na świecie zbożny pobyt} & \sig{K} c$'$ a c$'$ G c$'$ a-G-F G-F-G a \quad \sig{K70} c$'$ b\,natural c$'$ d$'$ G G c$'$ a G G b\,flat a \quad \sig{W} a c$'$ d$'$ G G c$'$ a F G a a\\
2.9 & \pl{po żywocie rajski przebyt} & \sig{K} a c$'$ d$'$ G c$'$ a-G-F G-F-G a (from \pl{raj-} on, identical to 2.8) \quad \sig{K70} a b\,flat c$'$ d$'$ G \textbar\ c$'$ a F F b\,flat a \quad \sig{W} a c$'$ d$'$ G c$'$ a F G a a\\
2.10 & \pl{Kyrieleison} & \sig{K} a G F F-G-a-G E D D \quad \sig{K70}, \sig{W} ``of no significance'' (Chybiński, no readings)\\
\bottomrule
\end{longtable}

\sect{Textual apparatus (stanzas 1--2)}

{\small Lemma in normalised form. Variants are normalised unless the spelling
matters; purely orthographic differences are ignored. Sigla as above.\par}

\begin{longtable}{@{}p{0.9cm}>{\raggedright\arraybackslash}p{3.3cm}>{\raggedright\arraybackslash}p{11.5cm}@{}}
\toprule
\textbf{Loc.} & \textbf{Lemma} & \textbf{Variants}\\
\midrule\endhead
1.1 & \pl{Bogu rodzica} & \sig{K} \sig{K408} \quad \pl{Boga rodzica} \sig{Wa} \sig{Cz} \sig{K3} \sig{Ł} \sig{M} \sig{H} \sig{S}\\
1.1 & \pl{sławiena} & \sig{K} \sig{K408} \quad \pl{sławiona} \sig{Wa} \sig{Cz} \sig{K3} \sig{Ł} \sig{M} \sig{H} \sig{S}\\
1.2 & \pl{U twego} & \sig{K} \sig{Wa} \sig{Cz} \sig{K3} \sig{Ł} \sig{H} \sig{S} \quad \pl{W twego} \sig{K408} \sig{M} \quad \pl{Twego} Woronczak's reconstructed text (editorial)\\
1.2 & \pl{Gospodzina} & \sig{K} \sig{K408} \sig{Wa} \sig{Cz} \sig{Ł} \sig{M} \quad \pl{gospodina} \sig{K3} \quad \pl{Gospodynia} \sig{H} \quad \pl{hospodyna} \sig{S}\\
1.2 & \pl{zwolena} & \sig{K} \sig{K408} \quad \pl{zwolona} \sig{Wa} \sig{Ł} \sig{M} \sig{H} \sig{S} \quad \pl{zbolenia} \sig{Cz} \quad \pl{zwolienia} \sig{K3}\\
1.3 & \pl{Zyszczy \dots\ spuści nam} & \sig{K} \sig{K408} \sig{Wa} \sig{Ł} \quad \pl{zyszczysz \dots\ spuszczysz} \sig{Cz} \quad \pl{zyszczy \dots\ spust winam} \sig{K3} \sig{M} \quad \pl{ziścisz \dots\ spuścisz} \sig{H} \quad \pl{ziści \dots\ spuści} \sig{S}\\
\midrule
2.1 & \pl{Twego dziela} & \sig{K} \sig{K408} \sig{Wa} \sig{Cz} \quad \pl{Twego syna} \sig{K3} \sig{Ł} \sig{M} \sig{H} \sig{S}\\
2.1 & \pl{Krzciciela} & all except \quad \pl{zbawiciela} \sig{M} \quad \pl{[Kr]ziciela zbawiciela} \sig{K3}\\
2.1 & \pl{Bożyce} & \sig{K} (\dipl{boszicze}) \sig{K408} (\dipl{bozide} Pilat) \quad \pl{zbożnica} \sig{Wa} \sig{Cz} \quad \pl{zbożny czas} \sig{Ł} \sig{H} \sig{S} \quad \pl{zbożnika} \sig{M}\\
2.2 & \pl{napełń} & \sig{K} \quad \pl{napełni} \sig{K408} \sig{Wa} \sig{Cz} \sig{K3} \sig{Ł} \sig{M} \sig{H} \sig{S}\\
2.3 & \pl{jąż nosimy} & \sig{K} \sig{K408} \sig{Wa} \quad \pl{jenżeć nosimy} \sig{Cz} \quad \pl{jenże (cię) prosimy} \sig{Ł} \sig{M} \sig{H} \sig{S}\\
2.4 & \pl{A dać raczy} & \sig{K} (Łoś, Kalina, Woronczak; \pl{O dać} Pilat) \sig{Wa} \quad \pl{Oddać raczy} \sig{K408} \sig{M} \quad \pl{O dać} \sig{Cz} \quad \pl{O o dać} \sig{Ł} \quad \pl{O o daci} \sig{H} \quad \pl{To dać} \sig{S}\\
2.5 & \pl{a na świecie} & \sig{K} \sig{K408} \quad \pl{daj na świecie} \sig{Wa} \sig{Cz} \sig{Ł} \sig{H} \sig{S} \quad \pl{daj nam na świecie} \sig{M} \quad \pl{dai na świecie dobry} \sig{K3}\\
2.5--6 & \pl{pobyt \dots\ przebyt} & \pl{pobytk \dots\ przybytk} \sig{Wa} (Woronczak; Pilat reads \dipl{przebythk})\\
2.6 & \pl{rajski} & \dipl{raski} \sig{K}\\
refr. & \pl{Kirielejson} & \dipl{kyrieleyson} \sig{K} (st.\,2 abbreviated \dipl{ky'eleyson}) \quad \dipl{kyon} \sig{Cz} \quad \dipl{kyeon} \sig{Wa} \quad \pl{Kyrie eleison} \sig{M} \sig{H} \sig{S}\\
\bottomrule
\end{longtable}

\subsection*{Notes}
\begin{description}[leftmargin=1.2em,itemsep=0.25em,font=\normalfont\itshape]
\item[1.1] \dipl{Bogv rodzicza} (\sig{K}, \sig{K408}) is two words: the
  dative \pl{Bogu} is used in a genitive sense (Łoś). Pilat preferred
  \pl{Boga rodzica}, the form in the 1410 \emph{Chronica conflictus} and in
  Długosz. \pl{Sławiena} and \pl{zwolena} are Old Polish forms without the
  \textit{e}\,>\,\textit{o} change, not necessarily Czechisms (Łoś, Kalina).
\item[1.2] \pl{U}: \sig{K} has two separate puncta d$'$ c$'$ here, which
  confirms a syllable before \pl{twe-} (Chybiński). The later musical
  witnesses \sig{K70} and \sig{W} lack both the syllable's notes and the
  octave leap.
\item[1.3] Kalina takes \pl{spust winam} in \sig{K3} and \sig{M} as proof of
  a shared exemplar for those two texts.
\item[2.1] \pl{dziela} (= \pl{dla}, `for the sake of'). Łoś saw this
  archaism as the main linguistic argument for a 13th-century date. Its
  replacement by \pl{syna} in the printed tradition, and by \pl{dzieła} in
  some modern editions, are both misunderstandings of it. \pl{Bożyce} is the
  vocative of \pl{bożyc}, `son of God'.
\item[2.4] On the decorated initial of \sig{K} see the palaeographical notes.
\end{description}

% ============================================================================
\sect{Palaeographical notes on \sig{K}}

These notes come from my collation of the Polona facsimile, 3293\,$\times$\,4699\,px.

\begin{itemize}[itemsep=0.25em]
\item \textbf{Staves and clefs.} System~1 has four ruled lines above the text:
  a C clef on the top line and an F clef on the second line from the bottom.
  The later systems have five ruled lines (Chybiński). The lowest of these
  also serves as the text baseline, which is why Kalina called the notation
  four-line. In system~2 the clef moves: after \pl{Gospodzina} a C clef is
  drawn on the middle line for \pl{Matko}, and before \pl{zwolena} an F clef
  restores the original position. Chybiński noted that the clefs are clearest
  ``at the beginning of the third and a third of the way from the end of the
  second system''. Read with these clefs, \pl{Matko} gives c$'$ a-G-F, as in
  every transcription.
\item \textbf{Faint or invisible notes.} The C on \pl{-gu-} (1.1) is very
  faded but traceable just above the text. The C on \pl{-pełń} (2.2) cannot be
  seen; only one punctum (D) stands over \dipl{naplen}. The first a of the
  final \pl{Kyrie} is faint. The ink blots over \pl{dziela} (system~4) and
  after \pl{a} (system~6) are stains, not notes. Everywhere else the reading
  is secure.
\item \textbf{Neume groups.} Over \pl{-dzi-na} of \pl{Gospodzina} the a and G
  are written close together and the F stands apart. This supports the
  underlay c$'$ a-G \,\textbar\, F rather than a single climacus.
\item \textbf{The initial of \pl{A dać raczy}.} It has the same design as the
  \pl{U} of \pl{Usłysz} two systems earlier: a dotted bowl on the left and a
  hatched box on the right. It is unlike the \pl{U} of \pl{U twego}. This may
  explain why Pilat read \dipl{O dacz}, while Kalina, Łoś and Woronczak read
  \dipl{A}. The later tradition splits the same way (\pl{O dać} \sig{Cz};
  \pl{O o dać} \sig{Ł}, \sig{H}). The notes are unaffected.
\item \textbf{The later hand.} Chybiński saw a cruder neume hand from \pl{A dać
  raczy} onwards. In the facsimile the puncta in systems~6--7 are larger and
  more widely spaced, but the forms are the same, and the pitches are clear
  once the staff slant is allowed for.
\end{itemize}

% ============================================================================
\sect{Commentary}

\textbf{Form.} Each stanza has three parts, two parallel lines followed by a
contrasting third and the refrain (Łoś). The melody is built by varying a small
stock of phrases. Surzyński (1891) labelled four of them \emph{a--d}. Pikulik
(2005) shows that Surzyński's fifth, \pl{Kyrieleison} motif is a variant of
\emph{d}. Flotzinger (1997/2005) set the phrases out in parallel. In his
account the whole melody can be traced back to the first line, which is built
almost symmetrically around a high point c$'$--d$'$--c$'$. Perz (2005) called
the procedure \emph{centonate}. The edition shows the main correspondences: the
cadence \pl{Maryja}\,=\,\pl{zwolena}\,=\,\pl{Maryja}\,=\,\pl{Bożyce}
(F-G-a-G E-C D); \pl{dziewica}\,=\,\pl{zyszczy nam}\,=\,\pl{usłysz
głosy}\,=\,\pl{a dać raczy}; \pl{Twego dziela Krzciciela}\,=\,\pl{napełń
myśli człowiecze}; and the repeat of the second half of \pl{a na świecie zbożny pobyt} in
\pl{po żywocie rajski przebyt}.

\textbf{Mode.} The final is D throughout, with cadences on a and F. Chybiński
weighed the hypolydian (sixth) mode, final F, because of the F-cadences on
\pl{sławiena}, \pl{Gospodzina}, \pl{Matko}, \pl{Kyrie}, and the lower fourth
C--F, but concluded for mode~1. B is avoided except once, on \pl{spuści}, at
the melodic peak.

\textbf{The octave leap.} \pl{Maryja} ends on D and \pl{U twego} begins on
d$'$. Pikulik argues that this leap could not occur before the 12th, and
properly the 13th, century. The later witnesses \sig{K70} and \sig{W} drop the
opening d$'$--c$'$ and so avoid the leap. That is most simply read as later
smoothing.

\textbf{Problem loci in \sig{K}.} (1)~\pl{Gospodzina}: Chybiński stresses that
the three puncta over the last two syllables do \emph{not} form a climacus,
since the first two are written apart from the third. They are printed here as
separate puncta (\texttt{j/h/g}). The spacing in the facsimile supports the
underlay \pl{-dzi-} c$'$-a-G, \pl{-na} F. (2)~The change of neume hand at \pl{A dać raczy}
makes the end of stanza~2 the least secure part of the source; the
correspondence with stanza~1 is the main control there. (3)~\pl{Kyrieleison}:
the underlay of seven syllables is modern. If the refrain was sung in five
(\pl{Ky-rie-e-lej-son}), the melisma F-G-a-G belongs with \pl{-e-} and E D D
with \pl{-lej-son}.

\textbf{Origins and date: positions in the literature.}
\begin{itemize}[itemsep=0.15em]
\item \emph{Attribution to St Adalbert} (d.\,997) is first found in 15th- and
  16th-century sources, among them Łaski 1506. It is rejected by all modern
  scholarship.
\item \emph{Two melodies by different authors} (Kalina 1880, on
  the analysis of Fr Bażański). Kalina
  separated a ``bright, periodic'' first part (stanza~1) from a ``gloomier''
  second part that he took to be Eastern in character. He found no link with
  the Czech \emph{Hospodine, pomiluj ny} (Prague, Univ. Lib. 10 E 2, f.~20).
  Feicht's 1962 analysis superseded this.
\item \emph{Late 13th century}, on linguistic and metrical grounds (Łoś 1922).
  Chybiński and Feicht place it at the turn of the 13th and 14th centuries.
\item \emph{Kyrie trope}: Woronczak (1952) identified the opening phrase with
  the Holy Saturday litany \emph{Kyrie} in the Graz gradual MS 807. Ostrowska
  added a neumed Minden source of 1024--27. Flotzinger (2005) shows that the
  agreement covers only the first phrase. He adds that the Graz MS was in fact
  written at St Nicholas, Passau, c.\,1170, and rejects the trope
  interpretation.
\item \emph{13th-century cento} built from litany, sequence and
  trouvère-style turns (Pikulik 2005). Feicht (1962) found a ``striking similarity''
  between the opening phrase and Jehan de Braine's \emph{Par desous l'ombre
  d'un bois}, one of many parallels he assembled (Gołos 1962). He concluded
  that the melody is homogeneous and the work of a trained musician.
\item \emph{Late 14th century, south-eastern Poland, with Czech mediation}
  (Flotzinger 2005). This rests on the Deesis structure of the text and on the
  Slavonic-rite foundations at Prague (Emaus, 1346), Oleśnica (1380) and
  Kraków (1390).
\item \emph{Ruthenian origin} (Shchurat 1906; Kuzminskyi 2017). This is a
  minority view. It is listed for completeness and is not adopted here.
\end{itemize}

% ============================================================================
\sect{Bibliography}
{\small
\begin{description}[leftmargin=1.5em,itemsep=0.15em,font=\normalfont]
\item Chmiel, A. ``Uwagi archiwalno-paleograficzne nad pieśnią \emph{Bogurodzica} w rękopisie Biblioteki Jagiell. nr 1619.'' \emph{Rozprawy Wydziału Filologicznego AU} 40 (1904). [with facsimile; not seen]
\item Chybiński, A. \emph{``Bogurodzica'' pod względem historyczno-muzycznym}. Kraków, 1907 (offprint from \emph{Przegląd Powszechny} 94--95). RCIN, open access.
\item Dąbrówka, A. ``Matka pieśni polskich.'' \emph{Pamiętnik Literacki} 96/2 (2005).
\item Feicht, H. ``Wstęp muzykologiczny.'' In \emph{Bogurodzica}, ed. J. Woronczak. Wrocław, 1962. [not seen]
\item Feicht, H. \emph{Studia nad muzyką polskiego średniowiecza}. Opera musicologica 1. Kraków: PWM, 1975. [not seen; pp. 187--89 for the Latin version]
\item Flotzinger, R. ``Zur \emph{Bogurodzica}-Frage.'' \emph{Context of Musicology} 1 (Poznań, 1997): 15--19. Polish trans. ``Jeszcze o kwestii \emph{Bogurodzicy}.'' \emph{Pamiętnik Literacki} 96/2 (2005): 7--10.
\item Jachimecki, Z. \emph{Muzyka polska w rozwoju historycznym}, vol.~1/1. Kraków, 1948.
\item Adamiec, M., ed. \emph{Bogurodzica}. Wirtualna Biblioteka Literatury Polskiej, Uniwersytet Gdański. \url{https://literat.ug.edu.pl/bgrdzc/} (transcriptions after Woronczak 1962; facsimiles of BJ 1619, Łaski 1506 and the lost Częstochowa copy).
\item Gołos, J. [G. S. Golos]. Review of Woronczak et al., \emph{Bogurodzica}. \emph{The Polish Review} 7/4 (1962): 84--85; also \emph{Notes} 2nd ser. 20/1 (1962--63): 71--72.
\item Kalina, A. \emph{Rozbiór krytyczny pieśni ``Bogarodzica''}. Lwów: W. Spasowicz, 1880.
\item Korolko, M. \emph{Średniowieczna pieśń religijna polska}. Wrocław, 1980.
\item Kuzminskyi, I. ``History of the Origin of the \emph{Bogurodzica} Song.'' \emph{Lietuvos muzikologija} 18 (2017): 147--161.
\item Łoś, J., ed. \emph{Bogurodzica. Pierwszy polski hymn narodowy}. Lublin, 1922. (Wikisource transcription.)
\item Pilat, R. \emph{Pieśń ``Bogarodzica''. I. Restytucyja}. Kraków: Druk. Uniw. Jagiellońskiego, 1879.
\item Perz, M. ``Polskie `posłowie' do uwag Rudolfa Flotzingera.'' \emph{Pamiętnik Literacki} 96/2 (2005): 11.
\item Pikulik, J. ``Co melodia \emph{Bogurodzicy} mówi nam o czasie powstania pieśni.'' \emph{Pamiętnik Literacki} 96/2 (2005): 13--15.
\item Stäblein, B. ``Litanei.'' \emph{MGG} 8 (1960), cols. 993--95.
\item Surzyński, J. \emph{Polskie pieśni Kościoła katolickiego od najdawniejszych czasów do końca XVI stulecia}. Poznań, 1891.
\item Urbańczyk, S. ``Das altpolnische Lied \emph{Bogurodzica} im Vergleich mit der altkirchenslawischen und alttschechischen Lyrik.'' \emph{Annales Instituti Slavici} II/2 (Wiesbaden, 1970).
\item Woronczak, J. ``Tropy i sekwencje w literaturze polskiej do połowy XVI wieku.'' \emph{Pamiętnik Literacki} 43 (1952): 353--54.
\item Woronczak, J., ed.; Ostrowska, E. (linguistic introduction); Feicht, H. (musicological ed.). \emph{Bogurodzica}. Biblioteka Pisarzy Polskich A~1 (Liryka średniowieczna 1). Wrocław: Ossolineum, 1962. 416 pp. [standard critical edition, with facsimiles of all sources and full transcriptions of every musical version; seen only through Adamiec's transcriptions and the reviews]
\end{description}}

\end{document}
